\section{Initial-Plan Data and Certificate Construction}
\label{app:new-model}

\subsection{Queue-aware model and reference controllers}

The seven-state model in \eqref{eq:new-plant} uses vehicle counts normalized
by 100 vehicles and time measured in hours. In the primitive-demand order
$(11,12,10,22,21,20)$, the fixed internal demands are
\begin{equation}
\begin{aligned}
\boldsymbol d=(&74.85630,37.42815,18.714075,\\
              &44.85630,7.42815,0.748563).
\end{aligned}
\label{eq:new-incident-demands}
\end{equation}
The displayed decimals are rounded; the exact stored binary coefficients
define the implemented field and are included in every proof. With base
exterior admission 60 and $u^\star=(0.5,0.5,0.5,0.52)$, the regional
state $(30,30,15,15,5,5)$ is an equilibrium and the exterior queue is 2.
The total base offered input is 244.031538 normalized units/hour, of which
the exterior stream supplies 24.587\%. These demands are reconstructed from
the specified equilibrium and then remain fixed across all events and
reference choices. The separate stationary record is retained in the earlier-study research package.

To construct a feedback reference with regional totals $(N_1,N_2)$ and
reference exterior arrival $A$, let $v_i=v(N_i)$, set
$u_{20}^\star=0.52$, $z^\star=2$ and $\pi=1/2$. Steady class balances give
\begin{equation}
\begin{aligned}
q_{11}^\star&=d_{11}+d_{21}+\pi A,&
q_{22}^\star&=d_{22}+d_{12}+(1-\pi)A,\\
n_{11}^\star&=q_{11}^\star/v_1,&
n_{22}^\star&=q_{22}^\star/v_2,\\
u_{12}^\star&=\frac{d_{12}+d_{10}}{G(N_1)-q_{11}^\star},&
n_{10}^\star&=\frac{d_{10}}{u_{12}^\star v_1},\\
n_{20}^\star&=\frac{d_{10}+d_{20}}{0.52v_2},&
n_{21}^\star&=N_2-n_{22}^\star-n_{20}^\star,\\
u_{21}^\star&=\frac{d_{21}+\pi A}{v_2n_{21}^\star},&
u_{02}^\star&=A/120.
\end{aligned}
\label{eq:new-targets}
\end{equation}
Linearization uses the strict capacity branch at the reference. The
continuous-time LQR weights are
$Q=\operatorname{diag}(1,1,0.1,0.1,0.1,0.1,2)$ and $R=3000I_4$.
The installed negative feedback gains give
$u=u^\star+K(x-x^\star)$ for reference triples
$(N_1,N_2,A)=(20,22,72),(22,24,78),(24,26,84)$.
Certification uses the actual event forecast, accounts for its difference
from the controller-design arrival $A$, and includes the full equilibrium
residual computed with the stored coefficients.

The controller-design search also includes $(18,20)$, reference arrivals
$60,72,78,84$, and input penalties $3000,30000,300000,3000000$.
All 64 outcomes are retained. The three controllers are the lowest-cost candidates for three
distinct target pairs on the selection grid. Two additional 45-minute
events broaden the validation set.

\section{Enclosure Proof and Verified Implementation}
\label{app:new-enclosure-proof}

\subsection{Conservation coordinates and nonsmooth flow}

The model equations and destination definitions are specified in
Section~\ref{sec:new-problem}; \appref{app:new-model} supplies the controller-design
parameters. For admission fraction $\pi$, its physical incidence is
$b_a=(0,1,0,1-\pi,0,0,-1)^\top$. The transformation
\eqref{eq:new-conservation-coordinates} has inverse
\begin{equation}
 \begin{aligned}
 n_1&=y_1+y_2+y_3,&
 n_2&=y_4+y_5+y_6-y_7,\\
 n_{11}&=y_1,&n_{22}&=y_4-(1-\pi)y_7,\\
 n_{10}&=y_3,&n_{20}&=y_6,\qquad z=y_7.
 \end{aligned}
 \label{eq:new-inverse-coordinates}
\end{equation}
Direct multiplication gives \eqref{eq:new-conservation-incidence}. Adding
each receiving-class equation to its share of $\dot z=\Lambda-a$ cancels
admission, while $n_1+n_2+z=\sum_{i=1}^{6}y_i$. These identities use a fixed
destination fraction and explicit conservation of waiting vehicles.

The class production $q_{ij}=n_{ij}G(n_i)/n_i$ is evaluated through the
polynomial speed $G(n_i)/n_i$, extended continuously at zero. Under affine
feedback the two admission branches are affine in the physical state, and
their minimum is continuous and locally Lipschitz. Consequently the
closed-loop field is locally Lipschitz in state and measurable in time.
Bounded measurable residuals give unique absolutely continuous solutions
locally; the compact enclosures below continue them throughout each verified
finite plan.

At a branch surface, the generalized Jacobian is contained in the convex hull
of the smooth branch Jacobians. This follows from the definition as limits
of classical Jacobians at differentiability points and their convex hull.
Both branches must be bounded on any outer predictor set intersecting that
surface. Outward bounds that establish a strict branch gap can instead
restrict evaluation to the feasible branch. The inequalities
\eqref{eq:new-local-bounds} remain valid for every convex combination of the
retained Jacobians.

\subsection{Componentwise first-exit argument}

Consider one predictor interval and omit its index $l$. Let
$z_e=L(x-p)$ denote the transformed prediction error, distinct from the
physical queue $z$. Before an exit from $|z_e|<R$, the line segment from
$p(t)$ to $x(t)$ lies in the verified outer set. The Lipschitz mean-value
relation along this segment yields, almost everywhere,
\begin{equation}
 \dot z_e=A(t)z_e+\epsilon(t)+(Le_7)w(t),
 \label{eq:new-error-dynamics}
\end{equation}
where $A(t)$ belongs to the convex hull of transformed generalized Jacobians
on the segment, and
$\epsilon=L f_b(p,t)-L\dot p$ has $|\epsilon_i|\le\delta_i$.
For $v_i=|(z_e)_i|$, the diagonal contribution is bounded by $a_iv_i$;
the remaining terms are bounded by
$\sum_{j\ne i}b_{ij}R_j+\delta_i+|(Le_7)_i|\bar w_x=c_i(R)$.
Absolute continuity therefore gives the scalar differential inequality
\begin{equation}
 \dot v_i\le a_iv_i+c_i(R)
 \quad\text{almost everywhere before exit}.
 \label{eq:new-error-comparison-proof}
\end{equation}
At a zero error component the corresponding upper derivative has the same
cross-component and forcing bound. The integrating-factor comparison gives
$v_i(t_l+s)\le\phi_i(s)$ from
\eqref{eq:new-component-comparison}. Since
$\phi_i'(s)=e^{a_is}(a_ir_i+c_i(R))$, its derivative has constant sign.
Its maximum on $[0,h]$ is thus $\max\{r_i,\phi_i(h)\}$.
The strict inequality \eqref{eq:new-self-inclusion} contradicts a first
exit by continuity. Every component remains enclosed throughout the interval.

At the right endpoint the error around the old center is bounded by
$\phi(h)$. Replacing that center by $p_{l+1}(t_{l+1})$ adds
\begin{equation}
 \omega_l\ge
 |L(p_l(t_{l+1})-p_{l+1}(t_{l+1}))|.
 \label{eq:new-recentering-bound}
\end{equation}
The triangle inequality gives the propagated radius stated in the main
text. A controller or target change preserves the physical state $x$.
With the common transform $L$, the additional radius comes from predictor
recentering. If a verified implementation
uses distinct invertible transforms before and after a boundary, the valid
transfer is
\begin{equation}
 r_{\rm next}\ge
 |L_{\rm next}L_{\rm old}^{-1}|\phi(h)
 +|L_{\rm next}(p_{\rm old}-p_{\rm next})|.
 \label{eq:new-coordinate-transfer}
\end{equation}
All centers in this expression are evaluated at the same boundary time.
An initial physical box centered at $p(t_0)$ with radius $d_x$ is enclosed
in the transformed coordinates by $r_0\ge|L|d_x$.

For an affine physical or command row, write
$x=p+L^{-1}z_e$. Its largest value over $|z_e|\le R$ is
$h^\top p+|h^\top L^{-1}|R$. Maximizing the predictor term in time proves
\eqref{eq:new-affine-support}. At a service checkpoint the endpoint radius
is available, giving \eqref{eq:new-service-support}. Finite induction over
all intervals, branch changes, and command boundaries completes the proof
of Theorem~\ref{thm:new-plan-enclosure}.

\subsection{Verified numerical bounds and exact constant commands}

The implemented piecewise quartic field and quadratic predictor produce a
defect polynomial of degree at most eight on each affine-forecast interval.
The verifier collects coefficients in the conservation coordinates before
interval evaluation. Outward arithmetic bounds the polynomial ranges,
Jacobian entries, inverses where needed, and scalar comparison functions.
Propagation intervals are split at the exact times of forecast knots,
command-block boundaries, and service deadlines. Predictors and trial radii
may be generated numerically; plan acceptance requires the outward-rounded
trajectory bounds and constraint checks to pass.

For $a\ne0$, the comparison integral can be evaluated as
\begin{equation}
 \chi(a,h)=\frac{\operatorname{expm1}(ah)}{a},
 \qquad \chi(0,h)=h.
 \label{eq:new-stable-comparison}
\end{equation}
This is the same integral as \eqref{eq:new-component-comparison}.
The \texttt{expm1} evaluation avoids cancellation when the bound on a
queue-Jacobian diagonal entry is close to zero, including when its exact
value is zero. Regression checks compare positive, negative, zero,
near-zero, and subnormal rates with high-precision evaluation.

For $K_b=0$, actuator rows are constants. The installed binary command
coefficients are interpreted as exact binary rationals, and the gate
residuals are evaluated exactly:
\begin{equation}
 \begin{gathered}
 u_i,\quad 1-u_i\quad(i=1,\ldots,4),\\
 \tfrac34\pm(u_{21}-u_{12}),\qquad
 \tfrac34\pm(u_{20}-u_{02}).
 \end{gathered}
 \label{eq:new-exact-actuator-residuals}
\end{equation}
A nonnegative exact residual satisfies its original closed inequality;
an outward lower floating-point representation preserves this conclusion,
including exact zero. The verifier applies this rule only after checking
that the gain is identically zero. The twelve state, class, regional, and
storage support calculations remain unchanged. Regression checks
include exact boundary commands and one-unit-in-the-last-place violations
of box and paired-gate rows. The proof identifier includes both the command
coefficients and this exact support rule.

\subsection{Proof records, retained portions, and task coverage}

Each plan record binds the initial set, model and constraints, commands,
forecast and residual bounds, absolute times, service rows, checkpoint sets
and expiry. Checkpoints identify their parent proof. After rejection,
execution follows the corresponding remainder of that record; restricting
its universally verified trajectory family preserves containment and
pending service checks, as required by Theorem~\ref{thm:new-plan-closure}.

For clarity, a replacement accepted before a task's last deadline must
include every pending service time, even when it uses a different nominal
forecast or prediction horizon. Its actual disturbance must belong to that
replacement's specified family. Until replacement, the current proof's own
assumptions must remain valid. A forecast update, sensing error outside the
certified bound, or a changed receiving-supply model is part of a new
certification record; the execution theorem is applied with a valid record
at each activation. An installation set containing the actual state can be obtained by
propagating the active plan's bounds during computation and intersecting
that enclosure with a measurement set that also contains the actual state. Without an accepted replacement, the finite-horizon guarantee remains
valid through the active plan's original certified expiry.

\section{Structure-Preserving Reachable Sets}
\label{app:structure-proof}

\subsection{Independent positivity and queue information}

Recall $v(n)=G(n)/n$, extended polynomially at zero. The stored and derived
class balances in \eqref{eq:new-plant} give
\begin{equation}
 \begin{aligned}
 \dot n_{12}&=d_{12}-u_{12}v(n_1)n_{12},\\
 \dot n_{21}&=d_{21}+\pi a-u_{21}v(n_2)n_{21}.
 \end{aligned}
 \label{eq:structure-derived-classes}
\end{equation}
Together with the four stored-class equations and queue equation, each
zero-class face has nonnegative inflow whenever commands, primitive demands
and offered arrivals are nonnegative and $v\ge0$. For the installed
binary coefficients $(a,b,c)=(0.0014877,-0.29815,15.0912)$, exact arithmetic gives
$a>0$ and $4ac-b^2>0$. Completing the square proves
\begin{equation}
 v(n)=a\left(n+\frac{b}{2a}\right)^2+
                    \frac{4ac-b^2}{4a}>0.
 \label{eq:structure-positive-factor}
\end{equation}
The minimum is approximately $0.15314$ h$^{-1}$. This algebraic continuation
establishes positivity independently of the requested jam and storage bounds.
For a local solution, the sum of negative class parts has upper derivative
at most a local Lipschitz constant times that sum. Its zero initial value
therefore remains zero. Equation~\eqref{eq:new-conservation} then gives
$0\le M(t)\le M(t_0)+\int_{t_0}^t(d_1+d_2+\Lambda(s))\,ds$;
bounded arrivals exclude finite-time escape. The nine lower-state/class
rows follow from this independent positive-cone argument. Its conclusion
is boundary invariance; full-vector monotonicity is a separate property.

For constant gates the scalar queue comparison is order-preserving in its
initial value and arrival input. With constant extremal arrival $\lambda$
and $C=120u_{02}$, its branch solutions are
\begin{equation}
 \begin{aligned}
 z(s)&=\lambda/60+(z_0-\lambda/60)e^{-60s}
                    &&(60z\le C),\\
 z(s)&=z_0+(\lambda-C)s &&(60z\ge C).
 \end{aligned}
 \label{eq:structure-queue-flows}
\end{equation}
Interval exponential/logarithmic evaluation encloses branch crossings.
Affine forecast intervals use their outward extremal arrivals when the common
branch remains unresolved; exact forecast and command cutpoints are retained.
This supplies endpoint and continuous bounds for all measurable allowed
residuals. The queue-storage requirement is tested against their upper bound.
If this continuous enclosure proves a common branch and the forecast is
$\ell_0+\ell_1s$, sharper endpoint comparisons use
\begin{equation}
\begin{aligned}
 z(h)&=e^{-60h}z_0+(\ell_0\pm\bar w_x)\varphi_1+\ell_1\varphi_2
       &&\text{(supply)},\\
 z(h)&=z_0+(\ell_0\pm\bar w_x-C)h+\ell_1h^2/2
       &&\text{(gate)},\\
 \varphi_1&=(1-e^{-60h})/60,\qquad
 \varphi_2=(h-\varphi_1)/60.
\end{aligned}
\end{equation}
The minus and plus signs give the lower and upper solutions. Only the
endpoint interval is intersected with these bounds; the independently
established continuous enclosure remains in force.

An initial physical box must itself satisfy class/queue nonnegativity; this
is checked by exact supports of the class map. Subsequent verified boxes may
be intersected with independently established positivity and queue intervals,
then enclosed outward. This operation preserves the reachable family and
includes recentering error. The Jacobian bounds within a step must cover the entire segment between
the predictor and the true state. For the queue coordinate, use the convex
hull of the predictor queue range and the independently bounded continuous
true queue range, intersected with the original trial enclosure. This
restricted enclosure is the Jacobian corridor used below. Admission branches
can be excluded when this enclosure rules them out, including when the
predictor lies outside the true queue range.

\subsection{Signed propagation and robust variational remainder}

Consider one interval, omit $l$, and let $\mathcal J_b$ bound each retained
transformed branch Jacobian on that corridor. In
\eqref{eq:structure-affine-set}, $\xi$ is constant across the plan;
the boxed remainder may depend on it. Choose any fixed matrix $A^0$ and
differentiate the decomposition to obtain
\eqref{eq:structure-remainder-identity} exactly. For a fresh physical box
$x(t_0)\in\widehat x+[-d_x,d_x]$, initialize $p(t_0)=\widehat x$,
$\Gamma(0)=L\operatorname{diag}(d_x)$ and $\rho=0$ before numerical
conversion; generator conversion error is absorbed outward into $\rho$.
Precisely, the mean-value matrix satisfies
$L[f_b(x,t)-f_b(p,t)]=B(s)L(x-p)$, with $B(s)$ in the convex hull of
the retained transformed branch Jacobians along the predictor-to-state
segment. The corridor condition applies to the full physical error; the
remainder $\zeta$ enters through its algebraic decomposition.
Define a Metzler matrix
$\mathcal M$ and nonnegative forcing bound $b_\zeta$ by
\begin{equation}
 \begin{aligned}
 \mathcal M_{ii}&=\max_b\sup(\mathcal J_b)_{ii},\\
 \mathcal M_{ij}&=\max_b\sup|(\mathcal J_b)_{ij}|\quad(i\ne j),\\
 (b_\zeta)_i&\ge\max_b\sup_{J\in\mathcal J_b,\,F\in\mathcal F}
                  \sum_j|[(J-A^0)F]_{ij}|\\
             &\hspace{1em}+\delta_i+|(Le_7)_i|\bar w_x,
 \end{aligned}
 \label{eq:structure-metzler}
\end{equation}
where $\mathcal F$ contains $e^{A^0s}\Gamma(0)$ for all $s\in[0,h]$,
and $\delta_i\ge\sup_s|\epsilon_i(s)|$. Convexification over branches
preserves these bounds. Applying absolute-value comparison to the remainder
identity gives
\begin{equation}
 |\zeta(s)|\le e^{\mathcal Ms}\rho+Q(s)b_\zeta,
 \qquad Q(s)=\int_0^s e^{\mathcal Mv}\,dv.
 \label{eq:structure-remainder-bound}
\end{equation}
The exponential is nonnegative, and $Q(s)$ is componentwise increasing.
The exponential itself can decrease in some entries; its old-remainder
contribution is bounded over the entire interval. With outward vectors
\begin{equation}
 \begin{aligned}
 f_i&\ge\sup_{s\in[0,h]}\sum_j|(e^{A^0s}\Gamma(0))_{ij}|,\\
 r_i^\zeta&\ge\sup_{s\in[0,h]}(e^{\mathcal Ms}\rho)_i
                            +(Q(h)b_\zeta)_i,
 \end{aligned}
 \label{eq:structure-continuous-radii}
\end{equation}
the test $f+r^\zeta<R$ keeps every error component strictly inside the
trial enclosure. The Jacobian bounds are valid until a hypothetical first
exit, and the strict inequality rules out that exit. The positivity and
scalar queue bounds used here are established independently of this
containment test.

For any affine row $a_x^\top x\le b_x$, its continuous support is bounded by
\begin{equation}
 \sup_s a_x^\top p(s)+
 \sup_s\|a_x^\top L^{-1}e^{A^0s}\Gamma(0)\|_1+
 |a_x^\top L^{-1}|r^\zeta.
 \label{eq:structure-signed-output}
\end{equation}
At a scheduled endpoint, use the endpoint generator and remainder.
Exact constant-command rows and the independently
proved positivity/queue supports can establish their corresponding constraints
directly. Each row can use the tighter of its two valid bounds.

Let $\widehat\Gamma$ be the stored binary generator at the next endpoint,
and let $d_\Gamma$ enclose the row sums of
$|e^{A^0h}\Gamma(0)-\widehat\Gamma|$. With recentering bound $\omega$ from
\eqref{eq:new-recentering-bound}, a valid update for the remainder bound is
\begin{equation}
 \rho^+\ge d_\Gamma+e^{\mathcal Mh}\rho+Q(h)b_\zeta+\omega.
 \label{eq:structure-endpoint-update}
\end{equation}
It retains the same $\xi$ and includes all numerical conversion errors.
The scalar queue bound constrains the total physical queue error; its
generator and remainder summands can cancel, so neither is separately
contracted by that total-error bound. Finite induction proves
Proposition~\ref{prop:structure-plan}.

\subsection{Validated matrix functions and execution records}

For any matrix $A$, $\alpha\ge\|A\|_\infty$, $z=\alpha h$ and series order
$p$, the exponential and integral tails obey
\begin{equation}
 \begin{aligned}
 \left\|e^{Ah}-\sum_{k=0}^p\frac{A^kh^k}{k!}\right\|_\infty
       &\le e^z\frac{z^{p+1}}{(p+1)!},\\
 \left\|\int_0^h e^{As}\,ds-
          \sum_{k=0}^p\frac{A^kh^{k+1}}{(k+1)!}\right\|_\infty
       &\le he^z\frac{z^{p+1}}{(p+2)!}.
 \end{aligned}
 \label{eq:structure-matrix-tails}
\end{equation}
The second inequality follows by integrating the first pointwise tail and
using $e^{\alpha s}\le e^{\alpha h}$. The implementation uses $p=18$.
Signed coefficients $a_x^\top L^{-1}(A^0)^k\Gamma(0)$ are formed before
interval time evaluation. For a row $C_i$ and generator column $\Gamma_j$,
the transformed tail is bounded by
$\|C_i\|_1\,e^z z^{p+1}/(p+1)!\,\|\Gamma_j\|_\infty$.
Outward rounding covers every retained coefficient, tail, product and sum.

Rejection preserves the accepted plan's generator and remainder history,
structural bounds, command schedule and expiry. Its checkpoint box
$|L(x-p)|\le |\Gamma|\mathbf1+\rho$ contains the reachable states but
omits the correlations in that history. Continuing the existing proof uses
the history; a fresh proof can use the correlated set or a containing box,
provided it verifies the full chosen set. A timestamped observation box
contains the actual state under its measurement assumptions. The replacement
certificate covers that entire set and every pending deadline. These
continuation and replacement cases establish the induction in
Theorem~\ref{thm:new-plan-closure}.

\section{Network Allocation Conditions}
\label{app:network-assumptions}

\subsection{Allocation changes and branch structure}

Proposition~\ref{prop:network-localization} uses constant allocation shares.
For a known differentiable $P_a(t)$, differentiating $y_c=c+P_a(t)z$ gives
\begin{equation}
 \dot y_c=F(c,u,t)+P_a(t)\Lambda+\dot P_a(t)z.
 \label{eq:network-varying-allocation}
\end{equation}
With fixed reference coordinates $y_c=c+P_0z$ and actual shares $P_a(t)$,
the corresponding expression is
\begin{equation}
 \dot y_c=F(c,u,t)+P_0\Lambda+(P_a(t)-P_0)a.
 \label{eq:network-uncertain-allocation}
\end{equation}
The extra terms enter the predictor and uncertainty enclosure. At a specified
piecewise-constant share change, physical inventories remain continuous but
coordinates change. Propagate the same initial uncertainty symbol through
$\Gamma^+=L_a^+(L_a^-)^{-1}\Gamma^-$, and bound the remainder by
$|L_a^+(L_a^-)^{-1}|\rho^-$ plus outward conversion and predictor-recentering
error. The implemented
transfer study uses constant shares and retains one coordinate map throughout.

For fixed commands and independent queues, bounding each queue's diagonal
Jacobian entry separately gives valid comparison bounds for
\eqref{eq:network-jacobian}, independently of simultaneous admission-branch
combinations. The predictor-defect calculation enumerates combinations of
branches that the local tests have not excluded, with at most $2^r$
combinations. Shared receiving capacities, regional-state-dependent
admission, and feedback gates require Jacobian bounds that include their
actual couplings.

\section{Task-State Transitions and Temporal Consistency}
\label{app:task-state-compiler}

\subsection{State, operations and temporal consistency}

The canonical record $\theta_i$ contains an active regional priority and a
last-in, first-out stack of suspended priorities, each with region, label and
absolute expiry. Network operation is the default. Before a proposal at $t_i$,
$\mathsf{Age}_{t_i}$ removes an expired active task and pops expired predecessors until
an unexpired task is restored or the stack is empty. For a parsed operation
$r_i^A$,
\begin{equation}
 \theta_i^+=C(\mathsf{Age}_{t_i}\theta_i,r_i^A),\qquad
 (j_i,\Delta_i)=Q(\theta_i^+,t_i),
 \label{eq:task-state-compiler}
\end{equation}
where $j_i$ selects the soft objective. A regional task expiring at $e_i$
has $\Delta_i=e_i-t_i>0$; network operation has $\Delta_i=0$.
The same expiry rule applies between proposals.

\begin{table}[t]
\centering\small
\caption{Compiler transitions. Only the listed fields accompany the operation
name. Durations/extensions are integer minutes in $[1,120]$; nonempty labels
contain at most 80 characters.}
\label{tab:task-state-operations}
\begin{tabular}{@{}p{0.14\textwidth}p{0.22\textwidth}p{0.56\textwidth}@{}}
\toprule
Operation & Additional fields & Transition at time $t$\\
\midrule
Set & Region, duration, label & Clear the stack; install a task expiring at
$t+$ duration.\\
Interrupt & Region, duration, label & Save the active task with its original
expiry; install a temporary task expiring at $t+$ duration.\\
Continue & None & Preserve the surviving task and expiry.\\
Extend & Minutes & Add the stated time to the active task's expiry.\\
Resume & None & Restore the latest unexpired saved task; if none survives,
clear the interruption. With no stack or interruption, keep the current task.\\
Cancel; network & None & Clear active and suspended tasks; use network operation.\\
\bottomrule
\end{tabular}
\end{table}

The compiler rejects extra fields, extension of an absent task and
interruption depth above four. Invalid operations create no new planning
request. The canonical record derives from schema-valid model operations,
independently of evaluation labels and physical-plan admission. A valid operation may update the requested objective while the earlier
certified plan continues. The task record describes the requested objective;
the execution record identifies the plan currently in use. Installing a new
plan requires its own passing certificate.
The hard storage, gate and service requirements remain in the trusted contract.

\emph{Temporal invariant.} Each stored expiry equals creation time plus the
specified duration and explicit extensions. Set and interrupt establish this
identity; extend adds exactly its argument; continue, suspension and restoration
preserve it. Deletion removes a task without changing another expiry.
Induction over the operation table, combined with $\mathsf{Age}_t$, therefore leaves either
network operation or a task with positive remaining time. Agreement with the
operator's language is scored separately at the compiled-request output.

\section{Continuation, Repair Timing and Cost Guarantees}
\label{app:repair-cost}

\subsection{Certified continuation beyond plan expiry}
\label{sec:terminal-backup}

In this framework, continuation beyond finite-plan expiry requires a
certified terminal backup: every admitted plan must end in a set from which that backup
preserves the physical constraints indefinitely.

A terminal backup consists of a nonempty set $\mathcal K$ and an implementable
state-feedback policy $\mu_B$ for which a stored certificate establishes that
$\mathcal K$ is robustly invariant. Specifically, for every activation state
$x(s)\in\mathcal K$ and every uncertainty realization within the specified
bounds, the backup trajectory exists for all $t\ge s$ and satisfies the
authorized physical state and command constraints. Any initialization required by $\mu_B$ is included in the certified
implementation.

The following proposition gives a sufficient condition for constructing such
a backup.

\begin{proposition}
\label{prop:backup-lyapunov}
Suppose the implemented backup has closed-loop dynamics
\[
\dot x=g_B(x,\eta),\qquad \eta(t)\in\mathcal E,
\]
where $\mathcal E$ is compact, $\eta(\cdot)$ is measurable, and $g_B$ is
continuous and locally Lipschitz in $x$, uniformly over
$\eta\in\mathcal E$, on an open domain $O$. Let
$V_B:O\rightarrow\mathbb R_{\ge0}$ be continuously differentiable and define
\[
\mathcal K=\{x\in O:V_B(x)\le\ell_B\},
\]
where $\mathcal K$ is nonempty and compact. If, on a neighborhood of
$\mathcal K$,
\[
\nabla V_B(x)^\top g_B(x,\eta)
\le -\lambda_B V_B(x)+\bar d_B
\]
for all $\eta\in\mathcal E$, with
$\lambda_B>0$ and
$0\le\bar d_B\le\lambda_B\ell_B$, and if every state and command generated
on $\mathcal K$ satisfies the authorized physical constraints, then
$\mathcal K$ is robustly invariant under $\mu_B$.
\end{proposition}

\begin{proof}
Along every admissible backup trajectory,
\[
\frac{d}{dt}V_B(x(t))
\le -\lambda_B V_B(x(t))+\bar d_B
\]
almost everywhere while the trajectory remains in the stated neighborhood.
Scalar comparison gives
\[
V_B(x(t))
\le
e^{-\lambda_B(t-s)}V_B(x(s))
+
\frac{\bar d_B}{\lambda_B}
\bigl(1-e^{-\lambda_B(t-s)}\bigr).
\]
If $x(s)\in\mathcal K$, then
$V_B(x(s))\le\ell_B$ and
$\bar d_B/\lambda_B\le\ell_B$, so
$V_B(x(t))\le\ell_B$. The trajectory therefore cannot leave
$\mathcal K$. Because $\mathcal K$ is compact and contained in the domain of
the closed-loop dynamics, the solution extends for all future times. The
assumed state and command constraints then remain satisfied under the backup.
\end{proof}

When continuation beyond the finite horizon is required, a candidate plan may
additionally be required to terminate in the certified backup set. Let a plan
installed at time $t_i$ have expiry $E_i>t_i$, and let
$\mathcal R_i(t)$ denote its certified reachable-state enclosure. We call the
plan \emph{terminally certified} if it satisfies the usual physical and
service checks over $[t_i,E_i]$, every service requirement that remains pending at installation has its
deadline in $[t_i,E_i]$ and is certified at that deadline, and
\begin{equation}
\mathcal R_i(E_i)\subseteq\mathcal K.
\label{eq:backup-terminal-landing}
\end{equation}
The certificate also specifies an implementable transition from the finite
plan to $\mu_B$ at $E_i$, including any required controller initialization.

For the Lyapunov certificate of
Proposition~\ref{prop:backup-lyapunov}, a sufficient terminal test is
\begin{equation}
\sup_{x\in\mathcal R_i(E_i)}V_B(x)\le\ell_B.
\label{eq:backup-computable-landing}
\end{equation}
The planner may impose this terminal condition, and the checker verifies
the entire terminal reachable set. The following theorem shows that
terminal certification provides a continuation when no acceptable finite
replacement is available.

\begin{theorem}
\label{thm:backup-recursive-availability}
Suppose a certified terminal backup $(\mathcal K,\mu_B)$ is available and an
initial terminally certified plan is installed. At each decision time, a
replacement is installed only if it is terminally certified; otherwise, the
current verified plan continues, and its backup policy is activated at expiry
if no replacement has been installed. If the realized uncertainty remains
within the certified bounds, execution conforms to the active certificate, and
replacement times are locally finite, then a certified continuation exists
for all future time, the physical constraints remain satisfied, and every
authorized service requirement is satisfied at its scheduled deadline. These
guarantees hold even if all subsequent replacement proposals are rejected.
\end{theorem}

\begin{proof}
Consider first a terminally certified plan installed at time $t_i$. Its
finite-plan certificate guarantees the physical constraints and all pending
service requirements through its expiry $E_i$. By
\eqref{eq:backup-terminal-landing},
\[
x(E_i)\in\mathcal R_i(E_i)\subseteq\mathcal K.
\]
The terminal-backup certificate therefore provides an admissible continuation
from $E_i$ onward. Hence, the finite plan together with its backup provides a
certified continuation beyond the finite plan's expiry.

Now consider a decision time before expiry. If a proposed replacement is
rejected, the remaining portion of the current verified plan continues under
the same certificate, including its terminal connection to $\mathcal K$. If a
replacement is accepted, its own terminal certificate contains the actual
installation state, verifies the applicable physical and service requirements,
and establishes a new terminal connection to $\mathcal K$. Thus, either
decision leaves an active certified continuation.

Applying this argument successively over the locally finite sequence of
decision times preserves the certification property. If no further finite
replacement is obtained, the active plan eventually reaches its expiry in
$\mathcal K$, where the certified backup takes over.

Finally, every accepted replacement is required to preserve all still-pending
service requirements, while rejection preserves the corresponding checks of
the existing certificate. Hence each authorized service requirement remains
verified at its scheduled deadline.
\end{proof}

Continuation follows an accepted plan, the incumbent remainder, or the
terminal backup. A service requirement introduced during backup operation
receives execution authority only after a new certificate covers that
requirement and a valid terminal continuation.

The indefinite backup guarantee also requires an uncertainty description
under which sustained operation is physically possible. In particular, a
finite exterior queue cannot remain bounded if the admissible arrival class
permits persistent inflow above every achievable discharge rate. The existence
of a nonempty robustly invariant set $\mathcal K$ therefore requires arrival
and control bounds compatible with sustained operation. Any diversion,
upstream holding, or demand-management mechanism used to satisfy this
condition must be represented in the certified model.

The numerical studies in Section~\ref{sec:new-evidence} evaluate
finite-horizon certification through the recorded plan expiries. The result
above characterizes continuation beyond expiry when the specified system
includes a certified backup set and policy.

\subsection{Installation timing and implemented-grid ordering}

A search conducted during incumbent execution must certify installation
at $s>t$ from the reachable information at $s$, using the remaining absolute
schedule. Minimum intervention applies to the grid for that activation
record. The synchronous experiments record computation separately;
Section~\ref{sec:delayed-handover-results} tests preparation during
incumbent execution.

The checker certifies the command coefficients that will actually be
executed. Equation~\eqref{eq:repair-family} gives the interpolation-distance
formula in exact arithmetic. For stored finite-precision commands, the
implementation orders grid plans by the Euclidean distance of their
coefficient vectors from the original proposal, or verifies that the stored
grid already follows that order. The tested grid passes this ordering check,
and every installed plan has a certificate for its stored coefficients.

\subsection{A conditional finite-horizon performance bound}
Consider the same initial state and disturbance realization for two fixed
command plans on $[t,t+H]$. Suppose their solutions remain in a common
compact domain on which, uniformly in admissible uncertainty,
\[
\|f(x,u,w)-f(y,v,w)\|
\le L_x\|x-y\|+L_u\|u-v\|.
\]
For a fixed running cost $\ell$ and terminal cost $g$, assume
$|\ell(x,u)-\ell(y,v)|\le L_{\ell x}\|x-y\|+
L_{\ell u}\|u-v\|$ and $|g(x)-g(y)|\le L_g\|x-y\|$.
Let $J(P)=\int_t^{t+H}\ell(x_P,u_P)\,ds+g(x_P(t+H))$.

The following proposition bounds the possible change in traffic cost caused by repairing a plan.

\begin{proposition}
\label{prop:repair-cost}
Under these assumptions, for the Euclidean norm on the stacked fixed-block
command vector,
\begin{equation}
|J(P(\vartheta))-J(P^0)|
\le L_J\vartheta\|P^B-P^0\|_2,
\label{eq:repair-cost-bound}
\end{equation}
where
\begin{align*}
L_J&=L_{\ell x}L_u B_H+L_{\ell u}H+L_gL_u A_H,\\
A_H&=\begin{cases}(e^{L_xH}-1)/L_x,&L_x>0,\\H,&L_x=0,\end{cases}\\
B_H&=\begin{cases}(e^{L_xH}-1-L_xH)/L_x^2,&L_x>0,\\H^2/2,&L_x=0.\end{cases}
\end{align*}
\end{proposition}
\begin{proof}
The maximum instantaneous command difference is at most
$\delta=\vartheta\|P^B-P^0\|_2$. Gronwall's inequality gives
$\|x_{P(\vartheta)}(t+s)-x_{P^0}(t+s)\|\le L_u\delta A_s$.
Integrating the running-cost bound and adding the terminal bound yields
\eqref{eq:repair-cost-bound}, since $\int_0^H A_s\,ds=B_H$.
\end{proof}

The queue-aware MFD has locally Lipschitz polynomial class flows on the
specified model domain and a minimum of Lipschitz admission branches.
Bounded controls, a finite horizon and a compact common trajectory domain
therefore provide finite $L_x,L_u$; the linear TTS and fixed-priority rows
have finite cost constants. The domain must contain both trajectories,
including the rejected reference plan if it leaves the physical constraint
set. This condition is additional to certification of the repaired plan.
The bound is uniform over paired initial states and residuals if the same
constants cover them all; it also bounds their worst-case costs by the
standard inequality $|\sup J_1-\sup J_0|\le\sup|J_1-J_0|$.
The proposition applies to fixed plans with a common continuous cost model.
The ensemble reports realized costs for the specified policies and
priority weights. A numerical $L_J$ requires valid bounds on the model and
cost constants over a common domain containing both the repaired trajectory
and the rejected reference trajectory, including any part outside the
physical constraint set.

\FloatBarrier
